\documentclass[10pt,a4paper]{amsart}
\usepackage[margin=19mm]{geometry}
\usepackage{amsmath,amssymb,mathtools}
\usepackage[scr=rsfs,cal=euler]{mathalfa}
\usepackage{xfrac}
\usepackage[unicode,hidelinks,bookmarks=false]{hyperref}
\newcommand{\ie}{i.e.\ }
\newcommand{\cf}{cf.\ }
\newcommand{\resp}{respectively\ }
\newcommand{\Subsection}{\subsection}
\providecommand{\nobreakdash}{\nobreak\hbox{-}}
\setlength{\emergencystretch}{2em}
\multlinegap0pt
\newcommand{\F}{\mathcal{F}}
\newcommand{\Kc}{\mathcal{K}}
\newcommand{\R}{\mathbb R}
\newcommand{\ga}{\gamma}
\newcommand{\lan}{\langle}
\newcommand{\pd}{\partial}
\newcommand{\ran}{\rangle}
\newcommand{\sig}{\sigma}
\newcommand{\supp}{\operatorname{supp}}
\theoremstyle{plain}
\newtheorem{theorem}{Theorem}[section]
\newtheorem{lemma}[theorem]{Lemma}
\newtheorem{definition}[theorem]{Definition}
\newtheorem{proposition}[theorem]{Proposition}
\newtheorem{corollary}[theorem]{Corollary}
\newtheorem{claim}[theorem]{Claim}
\theoremstyle{remark}
\newtheorem{remark}[theorem]{Remark}
\title[Local well-posedness for the KdV equation on the half-line at the critical regularity $H^{-\frac34}$]{Local well-posedness for the KdV equation on the half-line at the critical regularity $H^{-\frac34}$}
\author{M\'arcio Cavalcante}
\date{Source: arXiv:2609.28009v1 (CC BY 4.0)}
\begin{document}
\maketitle

\noindent\emph{Source and licence.} The delivered work is the article shown in the title above, version arXiv:2609.28009v1 (CC BY 4.0), distributed under the Creative Commons Attribution 4.0 International licence, as recorded in the arXiv licence metadata for this exact version and shown on the abstract page saved with this item. The full licence notice, which carries the licence URL and the address of that abstract page, is the bundled file \texttt{licenses/arxiv-2609.28009-CC-BY-4.0.txt}.
\par\smallskip

\noindent\textit{Editorial scope.} Counted unit: the article's Lemma (source label lem:trace) (source0.tex lines 258--262, source label \texttt{lem:trace}): ``Let , , and let be defined as in with replaced by in . Then for every with , x (t) G(x, ) H ( t) C G N s , and is continuous into .''. It is delivered together with the article's complete original proof (source0.tex lines 264--320, one proof environment adjacent to the statement, 57 lines), copied line for line from the article's own text. Required context retained uncounted: lem:cutoff (lines 178--180); eq:X (lines 211--214); eq:Nnorm (lines 220--222); lem:KPV (lines 247--251). Editorial formatting only: an AMS \texttt{amsart} preamble that restates the article's own macro definitions and its own theorem declarations, and the theorem environments the retained lines use; the article's own class file is neither shipped nor input; the retained environments are renumbered sequentially in order of appearance, whereas the article prints them with its own numbering; the article's equation numbering is not reproduced and every cross-reference inside the retained text resolves to the wrapper's numbering. Exactly 2 bibliography entries are delivered, transcribed from the article's own thebibliography; All retained bodies are exact copies of the bundled \texttt{sources/211540/source0.tex}, so no omission receipt applies. No asset-inclusion command occurs inside any retained span and the package ships no image asset.


\section*{Retained context: lem:cutoff (source0.tex lines 178--180)}

\begin{lemma}\label{lem:cutoff}
Let $\alpha\ge0$ and $\theta\in C_c^\infty(\R)$. Then $\|\theta g\|_{H^\alpha(\R)}\le C\|\theta\|_{H^{\alpha+1}(\R)}\|g\|_{H^\alpha(\R)}$.
\end{lemma}


\section*{Retained context: eq:X (source0.tex lines 211--214)}

\begin{equation}\label{eq:X}
\|u\|_X := \| \mathcal{F}^{-1} [\mathbf{1}_{\mathbb{R}^2 \setminus D} \hat{u}] \|_{X^{-\frac{3}{4}, \frac{1}{2}, 1}} + \| \mathcal{F}^{-1} [\mathbf{1}_D \hat{u}] \|_{X^{-\frac{3}{4}, \frac{1}{2}}},\qquad
\|u\|_Y:=\|\lan\xi\ran^{-\frac34}\widehat u\|_{L^2_\xi L^1_\tau}.
\end{equation}


\section*{Retained context: eq:Nnorm (source0.tex lines 220--222)}

\begin{equation}\label{eq:Nnorm}
\|G\|_{\mathcal N}:=\big\|\F^{-1}\big[\lan\tau-\xi^3\ran^{-1}\widehat G\big]\big\|_X+\big\|\F^{-1}\big[\lan\tau-\xi^3\ran^{-1}\widehat G\big]\big\|_Y ;
\end{equation}


\section*{Retained context: lem:KPV (source0.tex lines 247--251)}

\begin{lemma}[\cite{KPV}, \cite{Holmer}]\label{lem:KPV}
Let $s\in\R$, $\ga=\frac{s+1}3\in[0,\frac12)$, and $\theta\in C_c^\infty(\R)$ with $\supp\theta\subset[-2,2]$. Then
$$\sup_{x\in\R}\|\theta(t)e^{-t\pd_x^3}\phi(x,\cdot)\|_{H^\ga(\R_t)}\le C\|\theta\|_{H^1(\R)}\|\phi\|_{H^s(\R)},$$
and $x\mapsto\theta(t)e^{-t\pd_x^3}\phi(x,\cdot)$ is continuous into $H^\ga(\R_t)$.
\end{lemma}


\section{The counted unit: Lemma (source label lem:trace) and its complete original proof}

\begin{lemma}\label{lem:trace}
Let $-1<s<\frac12$, $\ga=\frac{s+1}3$, and let $\|\cdot\|_{\mathcal N_s}$ be defined as in \eqref{eq:Nnorm} with $-\frac34$ replaced by $s$ in \eqref{eq:X}. Then for every $G$ with $\|G\|_{\mathcal N_s}<\infty$,
$$\sup_{x\in\R}\|\psi(t)\Kc G(x,\cdot)\|_{H^\ga(\R_t)}\le C\|G\|_{\mathcal N_s},$$
and $x\mapsto\psi(t)\Kc G(x,\cdot)$ is continuous into $H^\ga(\R_t)$.
\end{lemma}

\begin{proof}
Set $W:=\F^{-1}[\lan\sig\ran^{-1}\widehat G]$, so that $\|G\|_{\mathcal N_s}=\|W\|_{X}+\|W\|_{Y}$ (with $s$ in place of $-\frac34$). By density it suffices to prove the bound for $G\in\mathcal S(\R^2)$; the continuity in $x$ then follows from the uniform bound and the explicit formulas below. Let $\psi_0\in C_c^\infty(\R)$ with $\psi_0=1$ on $[-1,1]$ and $\supp\psi_0\subset[-2,2]$, applied to the variable $\sig$. Taking the Fourier transform in $x$,
$$\mathcal{F}_x {\Kc G}(\xi,t)=\frac{1}{2\pi}\, e^{it\xi^3}\int_\R\frac{e^{it\sig}-1}{i\sig}\,\widehat G(\xi,\tau)\,d\tau ,$$
and we decompose $\Kc G=I+II+III$ with
\begin{align*}
I&=\frac{1}{(2\pi)^2}\iint e^{ix\xi+it\xi^3}\frac{e^{it\sig}-1}{i\sig}\psi_0(\sig)\widehat G\,d\tau d\xi
=\sum_{k\ge1}\frac{i^{k-1}t^k}{k!}\,e^{-t\pd_x^3}\phi_k,\\ 
&\quad \quad \widehat{\phi_k}(\xi)=\frac1{2\pi}\int\sig^{k-1}\psi_0(\sig)\widehat G\,d\tau,\\
&II=\frac{1}{(2\pi)^2}\iint e^{ix\xi+it\tau}\frac{1-\psi_0(\sig)}{i\sig}\widehat G\,d\tau d\xi,\\
&III=-\,e^{-t\pd_x^3}\phi_3,\qquad \widehat{\phi_3}(\xi)=\frac1{2\pi}\int\frac{1-\psi_0(\sig)}{i\sig}\widehat G\,d\tau .
\end{align*}

\emph{Term $I$.} By Lemma \ref{lem:KPV} with $\theta_k=t^k\psi$, $\|\theta_k\|_{H^1}\le C^k$, we get $$\|\psi I(x,\cdot)\|_{H^\ga}\le\sum_k\frac{C^k}{k!}\|\phi_k\|_{H^s},$$ and since $|\sig^{k-1}\psi_0(\sig)|\le 2^{k-1}\lan\sig\ran^{-1}\cdot\sqrt5$,
$$\|\phi_k\|_{H^s}\le C2^{k}\big\|\lan\xi\ran^s\lan\sig\ran^{-1}\widehat G\big\|_{L^2_\xi L^1_\tau}=C2^k\|W\|_Y .$$

\emph{Term $III$.} Since $(1-\psi_0(\sig))/|\sig|\le\sqrt2\lan\sig\ran^{-1}$, $\|\phi_3\|_{H^s}\le C\|W\|_Y$, and Lemma \ref{lem:KPV} gives $\|\psi\, III(x,\cdot)\|_{H^\ga}\le C\|W\|_Y$.

\emph{Term $II$.} By Lemma \ref{lem:cutoff} it suffices to bound $\|II(x,\cdot)\|_{H^\ga}$. The time Fourier transform of $II(x,\cdot)$ is $(2\pi)^{-1}F_x(\tau)$, where
$$F_x(\tau):=\int_\R e^{ix\xi}\frac{1-\psi_0(\sig)}{i\sig}\widehat G(\xi,\tau)\,d\xi .$$
All bounds below are in terms of $|\widehat G|$, hence uniform in $x$; we drop $e^{ix\xi}$. Set
$$D_\tau:=\int_\R|\widehat G(\xi,\tau)|^2\lan\xi\ran^{2s}\lan\sig\ran^{-1}d\xi,\qquad \int_\R D_\tau\,d\tau\le 2\|W\|^2_{X^{s,\frac12}}\le2\|W\|_X^2 .$$

\emph{(i) $|\tau|\le8$.} By Cauchy--Schwarz, $|F_x(\tau)|^2\le D_\tau\int\lan\xi\ran^{-2s}\lan\sig\ran^{-1}(1-\psi_0(\sig))^2d\xi$. For $|\xi|\le4$ the last integrand is bounded, and for $|\xi|\ge4$ one has $|\sig|\ge|\xi|^3-8\ge\frac12|\xi|^3$, so the integral is $\le C\int_{|\xi|\ge4}|\xi|^{-2s-3}d\xi+C\le C$ because $s>-1$. Hence $\int_{|\tau|\le8}\lan\tau\ran^{2\ga}|F_x|^2d\tau\le C\|W\|_X^2$.

\emph{(ii) $|\tau|\ge8$.} Let $r=r(\tau)$ be the real cube root of $\tau$, so $|r|\ge2$, and split $\R_\xi=R_1\cup R_2\cup R_3$,
$$R_1=\{|\xi-r|\le\tfrac12|r|\},\qquad R_3=\{|\xi|\ge\tfrac32|r|\},\qquad R_2=\R\setminus(R_1\cup R_3).$$
Note that $\xi\in R_1$ means that $\xi$ has the sign of $r$ and $\frac12|r|\le|\xi|\le\frac32|r|$. Hence a point of $R_2$ either satisfies $|\xi|<\frac12|r|$, whence $|\xi^3|<\frac18|\tau|$ and $|\sig|=|\tau-\xi^3|\ge\frac78|\tau|$, or has the sign opposite to $r$ (otherwise, having $\frac12|r|\le|\xi|<\frac32|r|$, it would lie in $R_1$), whence $\tau$ and $\xi^3$ have opposite signs and $|\sig|=|\tau|+|\xi|^3\ge|\tau|$. In both cases $|\sig|\ge\frac78|\tau|$ on $R_2$. On $R_3$, $|\sig|\ge|\xi|^3-|\tau|\ge\frac{19}{27}|\xi|^3$. On $R_1$, $\frac12|r|\le|\xi|\le\frac32|r|$, hence $\lan\xi\ran\sim|r|=|\tau|^{1/3}$ and $|\pd_\xi\sig|=3\xi^2\ge\frac34r^2$.

On $R_2$, by Cauchy--Schwarz and $s<\frac12$,
$$|F_{x,R_2}(\tau)|^2\le D_\tau\,\frac{C}{|\tau|}\int_{|\xi|\le\frac32|r|}\lan\xi\ran^{-2s}d\xi\le C D_\tau|\tau|^{-1}|\tau|^{\frac{1-2s}3}=CD_\tau\lan\tau\ran^{-2\ga},$$
since $-1+\frac{1-2s}3=-\frac{2s+2}3=-2\ga$. On $R_3$, by Cauchy--Schwarz and $s>-1$,
$$|F_{x,R_3}(\tau)|^2\le D_\tau\, C\int_{|\xi|\ge\frac32|r|}|\xi|^{-2s-3}d\xi\le CD_\tau|r|^{-2s-2}=CD_\tau\lan\tau\ran^{-2\ga}.$$
Integrating in $\tau$, the contributions of $R_2$ and $R_3$ to $\|\lan\tau\ran^\ga F_x\|_{L^2}^2$ are $\le C\|W\|_X^2$.

On $R_1$ we use the Besov structure. For $k\ge0$ let $E_k(\tau)=\{\xi\in R_1:2^k\le\lan\sig\ran<2^{k+1}\}$, and let $|E_k(\tau)|$ denote its Lebesgue measure in $\xi$. We claim that
\begin{equation}\label{eq:Ek}
|E_k(\tau)|\le \tfrac{16}{3}\,2^{k}|\tau|^{-2/3}.
\end{equation}
Indeed, the condition $2^k\le\lan\sig\ran<2^{k+1}$ means $$\sig\in J_k:=\{\sig\in\R: 2^{2k}-1\le\sig^2<2^{2k+2}-1\},$$ which is the union of two intervals (one in $\sig>0$, one in $\sig<0$) each of length at most $2^{k+1}$, so $|J_k|\le2^{k+2}$. On $R_1$ the map $\xi\mapsto\sig(\xi)=\tau-\xi^3$ is strictly monotone ($\xi$ has the sign of $r$ and $|\xi|\ge\frac12|r|$), with $|\sig'(\xi)|=3\xi^2\ge\frac34r^2=\frac34|\tau|^{2/3}$. Hence, changing variables,
$$
|E_k(\tau)|=\int_{E_k(\tau)}d\xi=\int_{\sig(E_k(\tau))}\frac{d\sig}{|\sig'(\xi(\sig))|}\le\frac{|J_k|}{\frac34|\tau|^{2/3}}\le\frac{2^{k+2}}{\frac34|\tau|^{2/3}},
$$
which is \eqref{eq:Ek}. (Geometrically, $E_k(\tau)$ consists of two intervals on either side of $r$, at distance $\approx2^k/(3r^2)$ from $r$ and of width $\approx2^k/(3r^2)$: the shell $B_k$ cut at height $\tau$. It is empty when $2^k\gg|\tau|$, since $|\sig|\le|\tau|+|\xi|^3\le\frac{35}{8}|\tau|$ on $R_1$.) Moreover $\lan\xi\ran^{-2s}\le C|\tau|^{-2s/3}$ on $R_1$, because $\frac12|r|\le|\xi|\le\frac32|r|$ and $|r|\ge2$. Hence, by Cauchy--Schwarz and \eqref{eq:Ek},
\begin{equation}
\begin{split}
\Big|\int_{E_k(\tau)}\frac{1-\psi_0(\sig)}{i\sig}\widehat G\,d\xi\Big|&\le 2^{1-k}\|\lan\xi\ran^s\widehat G(\cdot,\tau)\|_{L^2(E_k(\tau))}\Big(\int_{E_k(\tau)}\lan\xi\ran^{-2s}d\xi\Big)^{1/2}
\\&\le C2^{-k/2}|\tau|^{-\frac{s+1}3}\|\lan\xi\ran^s\widehat G(\cdot,\tau)\|_{L^2(E_k(\tau))}.
\end{split}
\end{equation}
Let $j(\tau)$ be defined by $2^{j(\tau)}\le\lan r\ran<2^{j(\tau)+1}$; then $$E_k(\tau)\times\{\tau\}\subset\bigcup_{|j-j(\tau)|\le2}(A_j\cap B_k),$$ and since $|\xi|\ge1$ on $R_1$, these sets lie outside $D$. Therefore
$$\lan\tau\ran^\ga|F_{x,R_1}(\tau)|\le C\sum_{|j-j(\tau)|\le2}\ \sum_{k\ge0}2^{-k/2}\|\lan\xi\ran^s\widehat G(\cdot,\tau)\|_{L^2_\xi(A_j\cap B_k)} .$$
Fix $j_0\ge0$. Taking the $L^2_\tau$ norm over $\{\tau:|\tau|\ge8,\ j(\tau)=j_0\}$ and using Minkowski's inequality,
$$\big\|\lan\tau\ran^\ga F_{x,R_1}\big\|_{L^2(\{j(\tau)=j_0\})}\le C\sum_{|j-j_0|\le2}\sum_{k\ge0}2^{-k/2}\|\lan\xi\ran^s\widehat G\|_{L^2_{\tau,\xi}(A_j\cap B_k)}\le C\sum_{|j-j_0|\le2}a_j,$$
where $a_j:=\sum_k2^{k/2}\|\lan\xi\ran^s\mathbf 1_{\R^2\setminus D}\widehat W\|_{L^2(A_j\cap B_k)}$; here we used $\lan\sig\ran\le2^{k+1}$ on $B_k$, so that $2^{-k/2}|\widehat G|\le 2\cdot 2^{k/2}\lan\sig\ran^{-1}|\widehat G|=2\cdot2^{k/2}|\widehat W|$. Summing the squares over $j_0$,
$$\big\|\lan\tau\ran^\ga F_{x,R_1}\big\|_{L^2(|\tau|\ge8)}^2\le C\sum_{j_0}\Big(\sum_{|j-j_0|\le2}a_j\Big)^2\le C\sum_ja_j^2\le C\|W\|_X^2 .$$
Collecting (i) and (ii), $\|II(x,\cdot)\|_{H^\ga}\le C\|W\|_X$ uniformly in $x$, which completes the proof.
\end{proof}

\begin{thebibliography}{99}
\bibitem[Holmer]{Holmer} J. Holmer, {\it The initial-boundary value problem for the Korteweg-de Vries equation,} Comm. Partial Differential Equations 31 (2006), 1151--1190.
\bibitem[KPV]{KPV} C. E. Kenig, G. Ponce and L. Vega, {\it Oscillatory integrals and regularity of dispersive equations,} Indiana Univ. Math. J. 40 (1991), 33--69.
\end{thebibliography}
\end{document}
